\section{Menaces à la validité et limites}\label{sec:menaces}

\paragraph{Sur-ajustement au développement.} Toutes les variantes de l'agent ont été choisies sur les mêmes \NQDev{}
questions du développement. L'exemple du Plan (section~\ref{sec:recup}) montre qu'un gain sur ce jeu peut disparaître sur
le test. Les descriptions de codes ont en outre été écrites après avoir vu des échecs du développement.

\paragraph{Jeu de test consommé.} Le jeu de test a servi une fois pour v1.1, puis une seconde fois pour comparer les
variantes de planification. Il n'est plus un test indépendant pour l'agent ; c'est pourquoi les résultats de la boucle
restent non confirmés jusqu'à l'évaluation sur un jeu neuf.

\paragraph{Juge LLM et circularité.} Un LLM qui juge les réponses d'un LLM peut partager ses biais. Nous imposons un modèle
juge distinct du générateur, nous l'alimentons avec le texte des PDF plutôt qu'avec nos \term{chunks}, et nous prévoyons un
contrôle humain d'un échantillon (accord et kappa de Cohen). Le vérificateur de support, lui, utilise le même modèle que le
générateur ; son étalonnage par mutations mesure ce qu'il détecte, pas ce qu'il manquerait sur des erreurs plus subtiles.

\paragraph{Taille d'échantillon et regroupement.} Avec \NConceptsDev{} concepts de développement, une question déplace le
Hit@6 de 1/\NQDev{} et les intervalles par concept sont larges. Les versions d'une même question dans trois langues ne sont
pas indépendantes ; traiter les questions comme indépendantes surestimerait la précision.

\paragraph{Non-déterminisme.} Le contrôle de stabilité porte sur peu de questions, toutes en français et du même code
(section~\ref{sec:stabilite}). Les quotas de l'API (erreurs 429 réessayées) ralentissent les exécutions sans changer leur
résultat, mais une limite journalière peut interrompre une évaluation ; les réponses sont mises en cache pour la reprendre.

\paragraph{Qualité d'extraction des PDF.} Notre vérité de terrain repose sur une seconde extraction (pypdf) qui a ses propres
défauts : mots coupés, notes de bas de page en arabe illisibles, numérotation des sections différente dans un PDF regroupé.
L'audit (section~\ref{sec:corpus}) chiffre ces écarts plutôt que de les supposer nuls.

\paragraph{Langues.} Les codes sont en français ; les questions arabes et anglaises passent par une traduction. Les effectifs
par langue sont trop petits pour conclure qu'une langue est mieux servie qu'une autre. Aucun texte arabe n'est reproduit
dans ce rapport.
